\documentclass[11pt,a4paper]{article}
\usepackage[utf8]{inputenc}
\usepackage[T1]{fontenc}
\usepackage[ngerman]{babel}
\usepackage{lmodern}
\usepackage{amsmath,amssymb}
\usepackage{geometry}
\usepackage[most]{tcolorbox}
\usepackage{enumitem}
\usepackage{array}
\usepackage{xcolor}
\usepackage{tikz}
\geometry{left=16mm,right=16mm,top=15mm,bottom=17mm}
\setlength{\parindent}{0pt}
\setlength{\parskip}{4pt}
\renewcommand{\arraystretch}{1.25}

\definecolor{green}{HTML}{0F7D49}
\definecolor{greenlight}{HTML}{EEF8F1}
\definecolor{bluelight}{HTML}{EEF6FF}
\definecolor{orangelight}{HTML}{FFF5E6}
\definecolor{linegray}{HTML}{D7E2DA}
\definecolor{ink}{HTML}{14231A}

\newtcolorbox{arbeitsbox}[2][]{colback=#2,colframe=linegray,arc=3mm,boxrule=.7pt,left=2.5mm,right=2.5mm,top=2mm,bottom=2mm,#1}
\newcommand{\linespace}[1]{\vspace{#1}}
\newcommand{\schreiblinien}[1]{\foreach \i in {1,...,#1}{\par\noindent\rule{\linewidth}{0.35pt}\vspace{0.32cm}}}
\newcommand{\kurzlinien}[1]{\foreach \i in {1,...,#1}{\par\noindent\rule{0.95\linewidth}{0.35pt}\vspace{0.22cm}}}
\newcommand{\vektor}[3]{\begin{pmatrix}#1\\#2\\#3\end{pmatrix}}

\begin{document}
\pagestyle{empty}

{\Large\bfseries Arbeitsblatt: Schnittpunkt zweier Geraden im Raum}\hfill {\bfseries EF}\par
{\color{green}\rule{\linewidth}{1.2pt}}
\textbf{Name:} \rule{0.34\linewidth}{0.4pt}\hfill\textbf{Datum:} \rule{0.22\linewidth}{0.4pt}

\begin{arbeitsbox}{greenlight}
\textbf{Problem: Kopfball nach einer Ecke}\par
Brown schlägt eine Ecke in den Strafraum. Schlotterbeck läuft zum Ball. Im Modell wird Browns Flanke durch eine Gerade und Schlotterbecks Kopfbewegung durch eine zweite Gerade beschrieben. Gesucht ist der Punkt, an dem beide Geraden denselben Ortsvektor liefern. Nur wenn Ort und Zeit zusammenpassen, ist der Kopfball im Modell möglich.
\end{arbeitsbox}

\begin{arbeitsbox}{bluelight}
\textbf{Modellwerte}\par
\[
S_0=\vektor{14}{20}{1{,}7},\quad S_1=\vektor{11}{16}{2{,}0},\qquad
B_0=\vektor{0}{0}{0{,}3},\quad B_1=\vektor{4}{6}{1{,}3}.
\]
Die beiden Bewegungen werden mit demselben Parameter \(t\) modelliert. Ein Parameterschritt entspricht ungefähr \(0{,}7\,\mathrm{s}\).
\end{arbeitsbox}

\textbf{1. Geradengleichungen aufstellen}\par
Bestimme jeweils einen Richtungsvektor und ergänze die Geradengleichung in der Form
\[
g:\;\vec{x}=\vec{p}+t\cdot\vec{u}.
\]

\begin{arbeitsbox}{white}
\textbf{Schlotterbecks Laufweg}\par
\[
\overrightarrow{S_0S_1}=S_1-S_0=
\vektor{\hspace{0.7cm}}{\hspace{0.7cm}}{\hspace{0.7cm}}
\]
\[
g_S:\;\vec{x}=\vektor{\hspace{0.7cm}}{\hspace{0.7cm}}{\hspace{0.7cm}}+t\cdot
\vektor{\hspace{0.7cm}}{\hspace{0.7cm}}{\hspace{0.7cm}}
\]
\end{arbeitsbox}

\begin{arbeitsbox}{white}
\textbf{Browns Flanke}\par
\[
\overrightarrow{B_0B_1}=B_1-B_0=
\vektor{\hspace{0.7cm}}{\hspace{0.7cm}}{\hspace{0.7cm}}
\]
\[
g_B:\;\vec{x}=\vektor{\hspace{0.7cm}}{\hspace{0.7cm}}{\hspace{0.7cm}}+t\cdot
\vektor{\hspace{0.7cm}}{\hspace{0.7cm}}{\hspace{0.7cm}}
\]
\end{arbeitsbox}

\textbf{2. Gleichsetzen und in ein lineares Gleichungssystem überführen}\par
Setze die beiden Geradengleichungen gleich und schreibe anschließend die drei Koordinatengleichungen auf.

\begin{arbeitsbox}{orangelight}
\[
\vektor{14}{20}{1{,}7}+t\cdot\vektor{-3}{-4}{0{,}3}
=
\vektor{0}{0}{0{,}3}+t\cdot\vektor{4}{6}{1}
\]
\[
\begin{aligned}
X_1:\quad && \rule{0.72\linewidth}{0.35pt}\\[0.25cm]
X_2:\quad && \rule{0.72\linewidth}{0.35pt}\\[0.25cm]
X_3:\quad && \rule{0.72\linewidth}{0.35pt}
\end{aligned}
\]
\end{arbeitsbox}

\textbf{3. LGS systematisch lösen}\par
Wähle ein geeignetes Verfahren und notiere deine Rechenschritte. Nutze anschließend den gefundenen Parameter, um den Schnittpunkt zu berechnen.

\begin{arbeitsbox}{white}
\textbf{Mein Lösungsverfahren:} \rule{0.65\linewidth}{0.35pt}
\schreiblinien{5}
\[
t=\rule{1.7cm}{0.35pt}\qquad H=\vektor{\hspace{0.8cm}}{\hspace{0.8cm}}{\hspace{0.8cm}}
\]
\end{arbeitsbox}

\newpage

\textbf{4. Rechenbeispiel: Einsetzungsverfahren}\par
Gegeben sind
\[
g_E:\;\vec{x}=\vektor{1}{2}{0}+t\cdot\vektor{2}{1}{1},\qquad
h_E:\;\vec{x}=\vektor{2}{10}{5}+r\cdot\vektor{1}{-2}{-1}.
\]

\begin{arbeitsbox}{greenlight}
\textbf{Schrittfolge}\par
Nach dem Gleichsetzen erhält man:
\[
\begin{aligned}
1+2t&=2+r\
2+t&=10-2r\
t&=5-r
\end{aligned}
\]
Aus der dritten Gleichung folgt \(r=5-t\). Setze das in die erste Gleichung ein:
\[
1+2t=2+(5-t)=7-t.
\]
Damit gilt \(3t=6\), also \(t=2\). Danach ist \(r=3\) und
\[
P=\vec{x}_{g_E}(2)=\vektor{1}{2}{0}+2\cdot\vektor{2}{1}{1}=\vektor{5}{4}{2}.
\]
\end{arbeitsbox}

\textbf{5. Übung: Additionsverfahren}\par
Gegeben sind
\[
g_A:\;\vec{x}=\vektor{4}{-1}{1}+t\cdot\vektor{1}{3}{1},\qquad
h_A:\;\vec{x}=\vektor{12}{2}{0}+r\cdot\vektor{-2}{1}{1}.
\]

\begin{arbeitsbox}{white}
\textbf{a) Überführe den Ansatz in ein LGS.}\par
\[
\begin{aligned}
X_1:\quad && \rule{0.72\linewidth}{0.35pt}\\[0.25cm]
X_2:\quad && \rule{0.72\linewidth}{0.35pt}\\[0.25cm]
X_3:\quad && \rule{0.72\linewidth}{0.35pt}
\end{aligned}
\]
\textbf{b) Löse mit dem Additionsverfahren.}\par
\schreiblinien{6}
\[
t=\rule{1.5cm}{0.35pt}\qquad r=\rule{1.5cm}{0.35pt}\qquad Q=\vektor{\hspace{0.7cm}}{\hspace{0.7cm}}{\hspace{0.7cm}}
\]
\end{arbeitsbox}

\begin{arbeitsbox}{bluelight}
\textbf{Worauf du beim Lösen eines LGS achten solltest}\par
\begin{itemize}[leftmargin=*]
\item Gleiche immer die ganzen Ortsvektoren, nicht nur eine Koordinate.
\item Schreibe die Koordinatengleichungen sauber untereinander und behalte die Parameter bei.
\item Prüfe den gefundenen Parameterwert in allen drei Gleichungen.
\item Ein Schnittpunkt im Raum ist erst gesichert, wenn alle Koordinaten denselben Punkt bestätigen.
\end{itemize}
\end{arbeitsbox}

\textbf{Merksatz}\par
Ergänze die fehlenden Wörter: \emph{gleichsetzen, LGS, Parameter, Schnittpunkt, prüfen}.
\begin{arbeitsbox}{orangelight}
Um den \rule{3cm}{0.35pt} zweier Geraden im Raum zu bestimmen, muss man die Geradengleichungen zunächst \rule{3cm}{0.35pt}. Daraus entsteht ein \rule{2cm}{0.35pt}. Die gefundenen \rule{2.7cm}{0.35pt} müssen anschließend in allen Koordinaten \rule{2.5cm}{0.35pt} werden.
\end{arbeitsbox}

\end{document}
